\documentclass[../ap_2013.tex]{subfiles}

\graphicspath{{./image/}}

\begin{document}

\setcounter{chapter}{0}
\chapter{古典力学:剛体}
\section{}
\(\vb*{F}\)は円柱Bと円柱Cとの接点で、左下の傾斜\(\ang{30}\)方向を向いている。
\(\vb*{F}'\)は床と円柱Bの接点で左向きを向いている。

円柱Bが回らないことから、モーメントの和は0で
\begin{equation}\begin{aligned}[b]
    -a\times F+a\times F' &=0\\
    \abs{\vb*{F}} &= \abs{\vb*{F}'}
\end{aligned}\end{equation}
とわかる。

\section{}
円柱Cと円柱A,Bの間の垂直抗力を\(N\),
円柱A,Bと床の間との垂直抗力を\(N'\)とする。

系全体の鉛直方向の力の釣り合いより、
\begin{equation}\begin{aligned}[b]
    2N' &= 3mg &  \rightarrow N' &= \frac{3}{2}mg
\end{aligned}\end{equation}
円柱Cの鉛直方向の力の釣り合いより、
\begin{equation}\begin{aligned}[b]
    mg = 2\qty(\frac{\sqrt{3}}{2}N+\frac{1}{2}F)=\sqrt{3}N+F
\end{aligned}\end{equation}
円柱Bの水平方向の力の釣り合いより、
\begin{equation}\begin{aligned}[b]
    \frac{1}{N} = F' +\frac{\sqrt{3}}{2}F
\end{aligned}\end{equation}
これと、設問[1]で得られた\(F=F'\)を用いて、
未知数\(F,F',N\)を求めると、
\begin{equation}\begin{aligned}[b]
    F = F' &= \frac{2-\sqrt{3}}{2}mg\\
    N &= \frac{1}{2}mg
\end{aligned}\end{equation}
滑らず静止する条件は\(F<\mu N, F'<\mu' N'\)で表されるので、
\begin{equation}\begin{aligned}[b]
    F &< \mu N & \rightarrow \mu &> 2-\sqrt{3}\\
    F' &< \mu' N' &\rightarrow \mu' &> \frac{2-\sqrt{3}}{3}
\end{aligned}\end{equation}

\section{}
円筒Aの密度は
\begin{equation}\begin{aligned}[b]
    \rho = \frac{m}{\pi a^2l}
\end{aligned}\end{equation}
である。
回転中心を軸とする円筒座標を考えて慣性モーメントを積分で表して計算すると
\begin{equation}\begin{aligned}[b]
    I &= \int_0^l dz\int_0^{2\pi} d\theta \int_0^a rdr \rho r^2\\
    &= \frac{2m}{a}\int_0^a r^4\,dr = \frac{1}{2}ma^2
\end{aligned}\end{equation}

\section{}
Tでは並進運動をせず、回転運動のみをしている。
図2の\(v_0\)の方向の並進を正、\(\omega_0\)の方向の回転を正として、
ST間の距離を\(x\), Tでの回転速度を\(\omega\), Tに到達するまでの時刻を\(t\)とする。

並進運動について、摩擦力はずっと左向きにかかっていて、
運動方程式を位置で積分していくと
\begin{equation}\begin{aligned}[b]
    \dv{t}mv(t) &= -\mu mg\\
    \int \dv{x}{t}\dv{v}{x}dx &= -\mu g\int dx\\
    \frac{1}{2}v(t)^2-\frac{1}{2}v_0^2 &= -\mu gx
\end{aligned}\end{equation}
地点Tでは\(v(t)=0\)なので、
\begin{equation}\begin{aligned}[b]
    x &= \frac{v_0^2}{2\mu g}
\end{aligned}\end{equation}
またこの時刻は、運動方程式を時間で積分すると
\begin{equation}\begin{aligned}[b]
    v(t) -v_0 &= -\mu gt\\
    t &= \frac{v_0}{\mu g}
\end{aligned}\end{equation}
である。

次に回転運動の運動方程式を時間積分すると、
\begin{equation}\begin{aligned}[b]
    \dv{t}I\omega(t) &= -a\mu mg\\
    I\omega(t)-I\omega_0 &= -a\mu mgt\\
    \omega(t) &= \omega_0 -\frac{a \mu mg}{\frac{1}{2}ma^2}\frac{v_0}{\mu g} = \omega_0 -\frac{2v_0}{a}
\end{aligned}\end{equation}

\section{}
滑らずに転がり始めるのは並進速度\(v\)と円柱表面の回転速度\(a\omega\)が一致したときである。
並進速度と回転速度の正の向きの取り方に気を付けると、この条件は
\begin{equation}\begin{aligned}[b]
    v=-a\omega
\end{aligned}\end{equation}
と書ける。

これを満たす時について、運動方程式の時間時間積分したものを整理すると
\begin{equation}\begin{aligned}[b]
    &\begin{cases}
        mv(t) -mv_0 &= -\mu mgt\\
        I\omega(t) -I\omega_0 &= -a\mu mgt
    \end{cases}\\
    &\begin{cases}
        v(t) -v_0 &= -\mu gt\\
        \frac{1}{2}a\omega(t) -\frac{1}{2}a\omega_0 &= -\mu gt
    \end{cases}\\
    v(t)-v_0 &= -\frac{1}{2}v(t)-\frac{1}{2}a\omega_0\\
    v(t) &= \frac{2v_0-a\omega_0}{3}
\end{aligned}\end{equation}
となる。これが求める速度である。


\section{}
設問[5]で得られた速度が負になっているので、
\begin{equation}\begin{aligned}[b]
    2v_0-a\omega_0 &< 0 &\rightarrow
    \frac{2v_0}{a} &<\omega_0
\end{aligned}\end{equation}

(4.1)式を初速度を円柱が位置Tにあるときの\(v=0\), 終速度を設問[5]にすると、
位置Tから滑らず転がり始める地点までの距離\(x'\)が右辺に出る。
なので、
\begin{equation}\begin{aligned}[b]
    \qty(\frac{2v_0-a\omega_0}{3})^2&=2\mu gx'\\
    x' &= \frac{1}{2\mu g}\qty(\frac{2v_0-a\omega_0}{3})^2
\end{aligned}\end{equation}
これが\(x>x'\)となる条件を調べているので、
\begin{equation}\begin{aligned}[b]
    \frac{v_0^2}{2\mu g} &> \frac{1}{2\mu g}\qty(\frac{2v_0-a\omega_0}{3})^2\\
    v_0 &> \frac{2v_0-a\omega_0}{3} > -v_0\\
    -v_0 &< a\omega_0 < 5v_0
\end{aligned}\end{equation}

よってこれらを満たすには
\begin{equation}\begin{aligned}[b]
    \frac{2v_0}{a} < \omega_0 < \frac{5v_0}{a}
\end{aligned}\end{equation}
であればよい。


\section*{感想}
古典力学のベクトルパズル苦手だなぁ。いっぱい解くしかない。

エネルギー保存則を使って解けないかと試したけど、
剛体が滑って回転しているときには筋が悪いとわかった。
摩擦熱\(\mu mg x\)の\(x\)というのは剛体表面と床がどれだけ擦れたかを表す量なので、
滑っているときには剛体の重心の移動距離から求められず、
全部の運動が求まってから最後にエネルギー保存則が成り立ってるねという形でチェックするしかなさそう。

\end{document}
